\documentclass[10pt,a4paper]{amsart}
\usepackage[margin=19mm]{geometry}
\usepackage{amsmath,amssymb,mathtools}
\usepackage[scr=rsfs,cal=euler]{mathalfa}
\usepackage{xfrac}
\usepackage[unicode,hidelinks,bookmarks=false]{hyperref}
\newcommand{\ie}{i.e.\ }
\newcommand{\cf}{cf.\ }
\newcommand{\resp}{respectively\ }
\newcommand{\Subsection}{\subsection}
\providecommand{\nobreakdash}{\nobreak\hbox{-}}
\setlength{\emergencystretch}{2em}
\multlinegap0pt
\usepackage{mathtools}
\usepackage{amssymb}
\usepackage{xcolor}
\usepackage[shortlabels]{enumitem}
\newtheorem{lemma}{Lemma}
\renewcommand{\arraystretch}{1.2}
\DeclareMathOperator{\rank}{rank}
\title{On the pluricomplex Poisson kernel and the associated complex Monge--Amp\`{e}re equation}
\author{Xieping Wang}
\date{Source: arXiv:2609.15572v1 (CC BY 4.0)}
\begin{document}
\maketitle

\noindent\emph{Source and licence.} On the pluricomplex Poisson kernel and the associated complex Monge--Amp\`{e}re equation. Version arXiv:2609.15572v1 (CC BY 4.0), distributed by arXiv under the Creative Commons Attribution 4.0 International licence, http://creativecommons.org/licenses/by/4.0/, as recorded in the arXiv licence metadata for this exact version and shown on the abstract page https://arxiv.org/abs/2609.15572v1. Full licence notice: \texttt{licenses/arxiv-2609.15572-CC-BY-4.0.txt}.\par\smallskip

\noindent\textit{Editorial scope.} Counted unit: the article's lemma lem:unique-maximizer (source0.tex lines 522--528). The article's complete original proof of it is delivered with it (source0.tex lines 530--586, one \texttt{proof} environment of 57 lines). No mathematics is written, repaired or re-derived: the statement and the proof are the article's own lines, in the article's own notation, and no retained line is substituted or re-flowed.  No further source-local context is needed: every cross-reference of every retained line resolves inside the two delivered spans.  Nothing was omitted from any retained span, so the package carries no omission record.  No retained line contains a citation, so the wrapper carries no bibliography and no bibliography entry.  No retained line contains a figure environment, an image-inclusion command or drawing code, so no image is delivered and the package declares no asset dependency.  Editorial formatting only, in addition: an AMS \texttt{amsart} preamble that restates the article's own declarations and macros used by the retained lines, namely the environments \texttt{lemma} on separate counters, together with the standard packages those lines need.  The retained environments print in this document as Lemma 1 in order of appearance, whereas the article prints the same items with the numbering of its own full document.  The complete native source of the article as distributed in the arXiv e-print for this exact version is bundled unchanged as \texttt{sources/210428/source0.tex}.
\begin{lemma}\label{lem:unique-maximizer} For every $\lambda\in(0,\, 1)$ and $(x,\, y)\in D$, the supremum defining $f_{\lambda}(x,\,y)$ is attained at a unique point $t=t_{\lambda}(x,\,y)$. Moreover, this maximizer depends smoothly on $(x,\, y)$  so that $f_{\lambda}$ is
smooth, with
   $$
   {\rm grad}\, f_{\lambda}=c_{\lambda}\circ t_{\lambda} \bigl(e^{t_{\lambda}},\, -1\bigr)
   \qquad {\rm and } \qquad  \rank {\rm Hess}_{\mathbb R}f_{\lambda}\equiv 1.
   $$
\end{lemma}

\begin{proof}
Given $\lambda\in(0,\, 1)$ and $(x,\, y)\in D$, we write $c=c_{\lambda}$ and
   $$
   E(t)\coloneqq c(t)\bigl(e^{t}(1+x-t)-y\bigr).
   $$
A straightforward calculation shows
   $$
   F(t)\coloneqq \frac{E'(t)}{e^{t}c(t)}
   =\bigl(1-\kappa(t)\bigr)(x-t)-\kappa(t)\bigl(1-ye^{-t}\bigr),
   $$
where
   $$
   \kappa(t)\coloneqq -\frac{c'(t)}{c(t)}=\frac{\lambda e^{t}}{(1+e^{t})^{2}}.
   $$
Since $0<\kappa(t)\leq \lambda/4 <1/4$, it follows that
 \begin{equation}\label{eq:F-limits}
    \lim_{t\to-\infty}F(t)=+\infty  \qquad {\rm and } \qquad \lim_{t\to+\infty}F(t)=-\infty.
 \end{equation}

We next show that $F$ admits exactly one zero point. Once this is proved, it follows immediately that $E$ attains its global strict maximum at this point,
as desired. To this end, first observe that at each zero point of $F$,
  \begin{equation}\label{eq:critical-identity}
    x-t = \frac{\kappa(t)}{1-\kappa(t)}\bigl(1-ye^{-t}\bigr).
  \end{equation}
From this and the fact that $y>e^{x}$, it follows that $ye^{-t}>1$. (Otherwise $1-ye^{-t}\geq0$, and the preceding equality would give $x\geq t$, hence $ye^{-t}>e^{x-t}\geq1$, a contradiction.) Note also that
   $$
    \kappa'(t)-\kappa(t)\bigl(1-\kappa(t)\bigr)
    =
    -\frac{\lambda e^{2t}(2+2e^{t}-\lambda)}{(1+e^{t})^{4}}
    <0.
   $$
Now differentiating $F$ and using this fact together with \eqref{eq:critical-identity}, we conclude that at each zero point of $F$,
  $$
  F'(t)=\left(\frac{\kappa'(t)}{1-\kappa(t)}-\kappa(t) \right)\bigl(ye^{-t}-1\bigr)-1 <-1.
  $$
It follows from this and \eqref{eq:F-limits} that $F$ admits exactly one zero point, which we denote by $t=t_{\lambda}(x,\,y)$,  and since $F'$ does not vanish there, the resulting
function $D\ni (x,\,y)\mapsto t_{\lambda}(x,\,y)$ is smooth by the implicit function theorem.

The envelope theorem now gives
  $$
  {\rm grad}\, f_{\lambda}=c_{\lambda}\circ t_{\lambda} \bigl(e^{t_{\lambda}},\, -1\bigr).
  $$
Differentiating this identity with respect to $x$ and $y$, respectively, yields
\begin{equation*}\label{eq:hessian-f}\renewcommand{\arraystretch}{1.8}
 {\rm Hess}_{\mathbb R} f_{\lambda}=\frac{\partial t_{\lambda}}{\partial x}
   \begin{pmatrix}
    e^{t_{\lambda}} (c_{\lambda}+c_{\lambda}')\circ t_{\lambda}  & -c_{\lambda}'\circ t_{\lambda}\\
     -c_{\lambda}'\circ t_{\lambda}  &  \dfrac{e^{-t_{\lambda}}(c_{\lambda}'\circ t_{\lambda})^{2}}
    {(c_{\lambda}+c_{\lambda}')\circ t_{\lambda} }
\end{pmatrix}.
\end{equation*}
(Note that $c_{\lambda}(t)+c_{\lambda}'(t)\neq0$ for all $t\in\mathbb R$, as is easily verified.) Similarly, differentiating  \eqref{eq:critical-identity}---with $t_{\lambda}$ in place of $t$---with respect to $x$ yields
  $$
  \frac{\partial t_{\lambda}}{\partial x}\neq0 \quad {\rm on }\ \, D.
  $$
Hence $\rank {\rm Hess}_{\mathbb R} f_{\lambda}\equiv 1$, and the proof is complete.
\end{proof}

\end{document}
